\documentclass[conference]{IEEEtran}
\usepackage{cite}
\usepackage{amsmath,amssymb}
\usepackage{booktabs}
\usepackage{array}
\usepackage{multirow}
\usepackage{xcolor}
\usepackage{url}
\usepackage{microtype}
\usepackage{balance}
\setlength{\tabcolsep}{3pt}
\newcolumntype{L}[1]{>{\raggedright\arraybackslash}p{#1}}
\newcommand{\archaic}{\textsc{Archaic}}
\newcommand{\designonly}{\textsc{Design-Only}}

\begin{document}

\title{Did Ancient Humans Interbreed with Other Hominins?\\
\archaic{} for Calibrated Inference of Reticulate Ancestry}

\author{Four-Stage Research Workflow\\
\textit{Survey - Idea - Experiment Design - Author}\\
\designonly{} research artifact\\
Evolutionary genomics and ancient-DNA research design}

\maketitle

\begin{abstract}
Ancient genomic data show that hominin history is not represented adequately by a single branching tree. A draft Neandertal genome revealed excess sharing with present-day Eurasians relative to sub-Saharan African comparison samples, a Denisovan genome showed a distinct archaic contribution to present-day Melanesians, and Denisova 11 was a first-generation offspring of a Neanderthal mother and Denisovan father. These results answer the headline question affirmatively, but each supports a different claim. An allele-sharing statistic does not by itself identify an exact donor, date, pulse count, or social context; a direct hybrid proves a cross-lineage reproductive event but not a population-wide mixture proportion. This paper introduces \archaic{}, Ancient-genome Reconstruction of Cross-Hominin Admixture and Introgression Claims, a seven-stage framework that links sample authenticity, reference representation, competing demographic models, local ancestry, timing and direction, independent validation, and calibrated reporting. The design uses damage and contamination bounds, D and $f_4$ statistics, mixture-model comparison, tract chronology, demography-aware ancestral recombination graphs, simulations, and held-out evaluation. It separates direct-hybrid, gene-flow-supported, source-qualified/time-supported, and adaptive-retention claims. The report is \designonly{}: it presents no new fossil sampling, sequencing, admixture fractions, tract maps, dates, or functional results.
\end{abstract}

\begin{IEEEkeywords}
ancient DNA, archaic introgression, Neanderthals, Denisovans, admixture graphs, D statistic, ancestral recombination graph, hominin evolution
\end{IEEEkeywords}

\section{Introduction}

Did ancient humans interbreed with other hominins? At the genetic level, the answer is yes: multiple independent data types show that at least some lineages exchanged ancestry. Green \emph{et al.} sequenced a draft Neandertal genome and found that Neandertal variants were shared more often with present-day Eurasians than with the sub-Saharan African individuals considered in their comparison, a pattern consistent with gene flow into ancestors of non-Africans~\cite{green2010}. Reich \emph{et al.} sequenced DNA from a Denisova Cave phalanx, established the Denisovans as a population sister to Neanderthals, and reported a contribution to present-day Melanesian genomes~\cite{reich2010}. The most direct individual-level evidence is Denisova 11, whose genome indicated a Neanderthal mother and a Denisovan father~\cite{slon2018}.

These findings changed the problem from ``did any contact occur?'' to a set of specific inference questions. Which lineages exchanged genes? In which direction? Was a signal produced by a single pulse, several episodes, ancestral population structure, or a source not yet sampled? Which retained segments were neutral passengers, which were removed by selection, and which, if any, made a functional contribution? The answer depends on a comparison set, an age model, a genome representation, and a statistical model. A percentage of ``archaic ancestry'' without those conditions is not a complete estimand.

This paper introduces \archaic{}, a framework for auditable claims about cross-hominin gene flow. It treats DNA damage and contamination as part of causal inference, does not equate an allele-sharing statistic with a detailed demographic history, and reserves functional claims for a separate evidentiary ladder. The aim is to turn a compelling historical conclusion into a reproducible research design rather than to add a new reconstruction.

\section{Evidence Base and Claim Boundaries}

\subsection{What is already well supported}

The Neandertal result is not based on morphology alone. The ancient genome comparison used genome-wide allele sharing and showed that a simple population tree could not explain all observed sharing patterns~\cite{green2010}. Subsequent high-coverage archaic genomes made it possible to study the genomic distribution of retained ancestry with greater resolution~\cite{prufer2014,meyer2012}. In present-day genomes, local ancestry and allele-sharing approaches recover Neandertal-derived haplotypes, while the distribution of those haplotypes indicates that selection and demographic history both shape what is retained~\cite{vernot2014,sankararaman2014,juric2016}.

Denisovan ancestry is more geographically and source heterogeneous than a single global label suggests. The first Denisovan genome paper inferred a contribution to present-day Melanesians~\cite{reich2010}; later diverse genome data supported multiple Denisovan-related source populations contributing to present-day genomes~\cite{bergstrom2020}. The Tibetan EPAS1 case is an important demonstration that an introgressed, Denisovan-like haplotype can have a population-specific adaptive history, but it does not imply that every archaic segment is adaptive~\cite{huerta2014}. A correct report keeps introgressed origin, selection, molecular function, and organismal effect as separate hypotheses.

Directly dated ancient modern-human genomes add a different scale of evidence. The Oase individual carried a recent Neanderthal ancestor~\cite{fu2015}; Initial Upper Paleolithic individuals from Bacho Kiro also showed recent Neanderthal ancestry~\cite{iasi2021}. These cases reveal recent genealogical connections for particular people and times. They should not be converted automatically into estimates for all contemporaneous populations. Conversely, the Denisova 11 genome supports actual cross-lineage parentage even though a single individual cannot establish how frequently the two populations met~\cite{slon2018}.

\begin{table*}[t]
\caption{Claim levels for cross-hominin interbreeding and the evidence required before promotion.}
\label{tab:claims}
\centering
\footnotesize
\begin{tabular}{L{0.18\textwidth} L{0.25\textwidth} L{0.28\textwidth} L{0.20\textwidth}}
\toprule
Claim level & What can be said & Minimum evidence & Main alternative or limitation \\
\midrule
Direct hybrid & A sampled individual has recent ancestry from two defined hominin populations & chromosome-scale ancestry pattern, authenticity controls, and simulated ILS/contamination comparisons & parental-source panel may be incomplete; individual result is not a population rate \\
Allele-sharing asymmetry & A declared quartet departs from a strict tree expectation & preregistered D or $f_4$ statistic, block-jackknife interval, and data-quality checks & ILS, ancestral structure, reference bias, or unmodeled gene flow \\
Gene-flow-supported & A gene-flow model predicts data better than named tree/structure alternatives & genome-wide statistics plus model comparison and sensitivity analysis & exact source, direction, time, and pulse count can remain unidentified \\
Source-qualified/time-supported & A source family and date interval are supported conditional on model assumptions & local ancestry/tract or serial ancient-genome evidence, recombination sensitivity, and competing source models & ghost lineage, multiple pulses, map uncertainty, and selection on tracts \\
Adaptive-retention candidate & An introgressed haplotype has evidence beyond introgressed origin & selection model, demographic controls, independent replication, and functional/phenotype evidence & linked selection, annotation bias, or cohort ascertainment \\
\bottomrule
\end{tabular}
\end{table*}

\subsection{What the headline question cannot settle}

The word ``interbreed'' can mean a documented first-generation child, a lineage-level gene-flow event, or merely shared ancestry. \archaic{} uses the first two meanings only when the genetic evidence supports them. It does not infer mating behavior, cultural acceptance, fertility of all descendants, or the social mechanisms of contact. Nor does a lineage name supply a unique genetic donor: a sampled Neanderthal or Denisovan genome is a reference, not necessarily the exact population that contributed ancestry.

Gene flow can also be reciprocal. Evidence for ancient-human-related input into eastern Neanderthals illustrates why a model with arrows in only one direction is a hypothesis, not a default historical fact~\cite{kuhlwilm2016}. Demography-aware ancestral-recombination-graph analysis can identify candidate older migration events, but source identity becomes especially uncertain as the divergence becomes deep or a donor is unsampled~\cite{hubisz2020,schaefer2021}. These limitations motivate a framework that reports a source family and explicit alternatives rather than a falsely precise fossil label.

\section{The \archaic{} Framework}

\archaic{} stands for Ancient-genome Reconstruction of Cross-Hominin Admixture and Introgression Claims. Its seven stages impose a sequence: establish that a sequence is interpretable before measuring allele sharing; establish an allele-sharing signal before assigning tracts; establish source/timing only after demographic alternatives; and establish function only after an introgression model has passed independent checks. The stage order guards against a common failure mode in which a visually appealing admixture graph is mistaken for unique history.

\begin{figure*}[t]
\centering
\fbox{\begin{minipage}{0.94\textwidth}
\centering\small
\begin{tabular}{c@{\quad$\Longrightarrow$\quad}c@{\quad$\Longrightarrow$\quad}c@{\quad$\Longrightarrow$\quad}c}
\fbox{\begin{minipage}{0.17\textwidth}\centering\textbf{A: Authenticity}\\Damage, contamination, age, coverage\end{minipage}} &
\fbox{\begin{minipage}{0.17\textwidth}\centering\textbf{R: Reference}\\Panels, masks, outgroups, maps\end{minipage}} &
\fbox{\begin{minipage}{0.17\textwidth}\centering\textbf{C: Competing models}\\Tree, structure, pulse, ghost source\end{minipage}} &
\fbox{\begin{minipage}{0.17\textwidth}\centering\textbf{H: Haplotypes}\\Local ancestry and genealogies\end{minipage}}\\\vspace{8pt}
\fbox{\begin{minipage}{0.17\textwidth}\centering\textbf{A: Admixture}\\Direction, timing, mixture interval\end{minipage}} &
\fbox{\begin{minipage}{0.17\textwidth}\centering\textbf{I: Independent validation}\\Replicate, holdout, simulation\end{minipage}} &
\fbox{\begin{minipage}{0.17\textwidth}\centering\textbf{C: Claim calibration}\\Direct, gene-flow, time/source, function\end{minipage}} &
\fbox{\begin{minipage}{0.17\textwidth}\centering\textbf{Versioned report}\\Evidence, alternatives, and limits\end{minipage}}
\end{tabular}
\end{minipage}}
\caption{\archaic{} converts a broad interbreeding question into linked evidence stages. Failure at a stage narrows the permitted claim; it does not erase a lower-level observation.}
\label{fig:workflow}
\end{figure*}

\subsection{A: authenticity and ascertainment}

For sample $j$, the study freezes a passport before population-genetic analysis:
\begin{equation}
\label{eq:passport}
\Pi_j=(S_j,\tau_j,B_j,C_j,\delta_j,\kappa_j,\Omega_j,\mathcal{R},\mathcal{O}),
\end{equation}
where $S_j$ is specimen and provenance information, $\tau_j$ the age model, $B_j$ the extraction/library batch, $C_j$ a contamination bound, $\delta_j$ the damage profile, $\kappa_j$ coverage and callability, $\Omega_j$ the included genomic regions, $\mathcal{R}$ the reference representation, and $\mathcal{O}$ the designated outgroups. Equation~\eqref{eq:passport} prevents a sample or mask from changing after a desired ancestry pattern is seen.

An ancient read can be modeled as a contamination mixture. For base $b$ at locus $\ell$, a generic likelihood is
\begin{equation}
\label{eq:authenticity}
\Pr(b_{j\ell}\mid g_{j\ell})=(1-c_j)\Pr(b_{j\ell}\mid g_{j\ell},\delta_j)+c_j\Pr(b_{j\ell}\mid g^{(c)}_{\ell}),
\end{equation}
where $g_{j\ell}$ is the endogenous genotype, $c_j$ is a contamination parameter, $\delta_j$ represents damage and sequencing error, and $g^{(c)}_{\ell}$ denotes the contaminant source distribution. Equation~\eqref{eq:authenticity} is a design template, not a claim that one contaminant model fits every fossil. The analysis reports results across credible $c_j$ values, read-end trimming choices, and transversion-only filters.

\begin{table*}[t]
\caption{ARCHAIC stages, required outputs, and fail-fast decisions.}
\label{tab:stages}
\centering
\footnotesize
\begin{tabular}{L{0.11\textwidth} L{0.23\textwidth} L{0.29\textwidth} L{0.27\textwidth}}
\toprule
Stage & Required output & Minimum controls & Fail-fast decision \\
\midrule
A: Authenticity & passport and usable data interval & damage, contamination, coverage, date/provenance, library/batch record & stop lineage inference if identity or quality is not bounded \\
R: Reference & panels, masks, outgroups, and map representation & reference-bias checks, callability and sex-chromosome policy & stop if labels or comparable genomic regions are undefined \\
C: Competing models & preregistered tree, structure, pulse, reciprocal, and ghost alternatives & simulated ILS and demographic sensitivity & do not promote a statistic to gene-flow claim without alternatives \\
H: Haplotypes & local ancestry/genealogy map with uncertainty & at least two methods or reference configurations & retain genome-wide signal only if tract assignment is unstable \\
A: Admixture & conditional direction, pulse, mixture, and date interval & recombination and source-panel sensitivity & downgrade exact time/donor language if intervals overlap alternatives \\
I: Independent validation & replicate, holdout, or independent data partition result & frozen specifications and disagreement analysis & mark claim provisional if it does not generalize \\
C: Calibration & evidence-labeled report and data/code provenance & claim ladder and deviation ledger & prohibit functional or social inference from ancestry alone \\
\bottomrule
\end{tabular}
\end{table*}

\subsection{R and C: reference representation and competing models}

Reference choice is a biological and computational decision. Samples are not reduced to an undifferentiated ``African,'' ``Eurasian,'' or ``archaic'' label; the report states individual, site, time-bin, and population labels separately. It records the reference build, mapping method, callable-site mask, outgroup, and whether observations are genotype likelihoods or pseudo-haploid calls. This matters because mapping and ascertainment can alter allele-sharing patterns, especially for damaged or low-coverage ancient sequences.

For quartet $(P_1,P_2,P_3,O)$, a standard test statistic is
\begin{equation}
\label{eq:dstat}
D=\frac{n_{\mathrm{ABBA}}-n_{\mathrm{BABA}}}{n_{\mathrm{ABBA}}+n_{\mathrm{BABA}}},
\end{equation}
where $P_1$ and $P_2$ are comparison populations, $P_3$ is a candidate archaic source, $O$ is an outgroup, and the numerator counts the difference between two discordant derived-allele configurations. Equation~\eqref{eq:dstat} is evaluated with genomic blocks and a block jackknife because sites are linked. Durand \emph{et al.} formalized testing of ancient admixture with such sharing patterns, while Patterson \emph{et al.} developed a broad framework for population-history tests~\cite{durand2011,patterson2012}.

The related $f_4$ statistic is written as
\begin{equation}
\label{eq:f4}
f_4(A,B;C,D)=\frac{1}{L}\sum_{\ell=1}^{L}(p_{A\ell}-p_{B\ell})(p_{C\ell}-p_{D\ell}),
\end{equation}
where $p_{X\ell}$ is the allele-frequency or uncertainty-aware allele-state estimate for group $X$ at locus $\ell$. Equation~\eqref{eq:f4} measures covariance in allele-frequency differences; its interpretation is conditional on the quartet and population history. A nonzero value motivates model comparison, not automatic identification of a single edge in a graph.

\section{Study Design for a New Analysis}

\subsection{Sampling architecture}

The first \archaic{} implementation uses a nested design. It includes: (i) authenticated ancient hominin references from more than one time or locality where possible; (ii) ancient modern-human genomes near relevant contact windows; (iii) geographically stratified present-day genomes under valid consent and access agreements; (iv) an outgroup appropriate to the comparison; and (v) genomic regions with a preregistered callability policy. A low-coverage fossil is not discarded by default; it is analyzed with its genotype uncertainty and is never allowed to silently drive a high-confidence local-ancestry conclusion.

Table~\ref{tab:sampling} makes the sampling commitments explicit. The required provenance and community governance are not administrative extras: destructive sampling and data access determine which comparisons are possible, and unrecorded exclusions can create a misleading population representation.

\begin{table*}[t]
\caption{Sampling and measurement architecture for a first ARCHAIC implementation.}
\label{tab:sampling}
\centering
\footnotesize
\begin{tabular}{L{0.18\textwidth} L{0.27\textwidth} L{0.27\textwidth} L{0.18\textwidth}}
\toprule
Evidence layer & Design element & Recorded controls & Primary use \\
\midrule
Ancient specimen & provenance, age model, extraction/library identity, endogenous DNA, coverage & permits, destructive-sampling record, damage, contamination, duplicate rate & authenticity and temporal anchoring \\
Archaic reference panel & multiple Neanderthal/Denisovan-related individuals where available & coverage class, callability, temporal/locality label, reference bias & source-family and genealogy comparison \\
Ancient modern human & dated individuals near contact windows & age uncertainty, kinship, population label, coverage and missingness & recent ancestry and serial validation \\
Present-day panel & stratified recipients and comparison populations & consent/access, sampling frame, relatedness, phasing/callability & population-level allele sharing and local ancestry \\
Genomic representation & fixed build, outgroup, mask, recombination maps & transversion subset, read position, autosome/sex-chromosome policy & robust statistics and timing sensitivity \\
Functional module & separately approved molecular/phenotype data when justified & annotation version, cohort ascertainment, demography and linked-selection controls & adaptive-retention assessment only \\
\bottomrule
\end{tabular}
\end{table*}

\subsection{From allele sharing to a mixture parameter}

If a source graph is specified, a mixture parameter can be represented with an $f_4$ ratio:
\begin{equation}
\label{eq:mixture}
\widehat{\alpha}=\frac{f_4(P_1,O;T,R)}{f_4(P_1,O;S,R)},
\end{equation}
where $T$ is a target, $S$ a proposed source, $R$ a reference related to the other ancestry component, $P_1$ a comparison population, and $O$ an outgroup. Equation~\eqref{eq:mixture} is meaningful only under its stated graph, source relation, and population labels. The study therefore reports $\widehat{\alpha}$ as a model-conditional interval, not a universal biological essence of a population.

Every graph is compared with alternatives. For data $Y$, model $M_k$, and parameters $\theta_k$, posterior or likelihood support has the generic form
\begin{equation}
\label{eq:modelsupport}
\Pr(M_k\mid Y)\propto \Pr(Y\mid M_k,\theta_k)\Pr(\theta_k\mid M_k)\Pr(M_k),
\end{equation}
where the likelihood incorporates sampling, error, linkage, and demographic assumptions. Equation~\eqref{eq:modelsupport} emphasizes that a preferred graph depends on the candidate set and priors or parameter constraints. A tree with ILS, a structured ancestral population, a single pulse, two pulses, reciprocal flow, and an unsampled-source model are all specified before viewing the primary results.

\subsection{Local ancestry and temporal inference}

Genome-wide asymmetry and local ancestry are complementary. A reference-based hidden Markov model can map segments that resemble an archaic panel; a demography-aware ancestral recombination graph can represent local tree variation and candidate migration under a user-defined model~\cite{hubisz2020,schaefer2021}. Agreement between methods is not a guarantee of truth, but disagreement is diagnostic: it can reveal weak source information, reference bias, recombination-map dependence, or an unmodeled population event.

Under a simplified single-pulse model, the expected tract length is approximately
\begin{equation}
\label{eq:tract}
\mathbb{E}[\ell]\approx\frac{1}{r t},
\end{equation}
where $\ell$ is inherited tract length, $r$ is a local recombination rate, and $t$ is time in generations since the pulse. Equation~\eqref{eq:tract} is not a universal dating equation. Selection, multiple pulses, source uncertainty, recombination variation, masking, and tract-calling errors all change the observed distribution. Dates must therefore be intervals assessed under alternative maps and demographic models. Recent early modern-human genomes can constrain timing in ways unavailable from present-day data alone~\cite{sumer2024}.

\begin{table*}[t]
\caption{Predictions that distinguish competing explanations of cross-hominin allele sharing.}
\label{tab:alternatives}
\centering
\footnotesize
\begin{tabular}{L{0.19\textwidth} L{0.23\textwidth} L{0.26\textwidth} L{0.22\textwidth}}
\toprule
Candidate explanation & Expected genome-wide pattern & Local/temporal prediction & Discriminating analysis \\
\midrule
Strict tree plus ILS & stochastic discordant sites without consistent source-recipient excess beyond the null model & short, diffuse genealogical variation without a coherent source-enriched tract set & coalescent simulations, D/$f_4$ block behavior, and ARG posterior comparison \\
Ancestral population structure & broad allele-sharing asymmetry can occur without post-split contact & signal need not follow a recent tract-length or dated-sample trajectory & structured-ancestor graph, serial ancient genomes, and source-panel sensitivity \\
Single pulse & coherent recipient-source sharing under one contact interval & tract length has a single-pulse-like distribution after accounting for maps and selection & tract likelihood, time-series genomes, held-out chromosome prediction \\
Multiple pulses or sources & heterogeneous ancestry levels and different source affinities across populations/segments & multimodal or population-specific tract/genealogy patterns & multi-source graph, local ancestry clustering, and diverse recipient panel \\
Reciprocal gene flow & asymmetry depends on time and selected reference orientation & introgressed segments can appear in archaic and modern-human-related branches & bidirectional graph and ancient reference chronology \\
Unsampled source & intermediate or inconsistent affinity to known archaic references & local trees may fit neither sampled source cleanly & ghost-source model, ARG inference, and conservative source-family language \\
\bottomrule
\end{tabular}
\end{table*}

\section{Evaluation and Claim Calibration}

\subsection{Direct-hybrid analysis}

For a specified ancient individual, a direct-hybrid claim requires more than an average ancestry proportion. It tests whether paternal and maternal chromosome-scale segments, heterozygosity, and ancestry switches are more probable under recent parentage than under population-level admixture, ILS, or contamination. A generic direct-hybrid comparison is
\begin{equation}
\label{eq:hybridlr}
\Lambda=\log\frac{\Pr(Y\mid H_{\mathrm{recent\ hybrid}},\Pi)}{\Pr(Y\mid H_{\mathrm{population\ admixture}},\Pi)},
\end{equation}
where $Y$ is the sequence evidence and $\Pi$ is the sample passport in \eqref{eq:passport}. Equation~\eqref{eq:hybridlr} is calibrated using damage-aware simulations with matched coverage, contamination, and source-panel gaps. Denisova 11 demonstrates why this individual-level category is distinct: the result supports a Neanderthal mother and Denisovan father for one sampled person~\cite{slon2018}, not a general frequency of such unions.

\subsection{Predictive performance, replication, and sensitivity}

Fitting an observed statistic is insufficient when several demographic graphs can fit it. With held-out unit $h$, the log predictive score is
\begin{equation}
\label{eq:lps}
\mathrm{LPS}(M_k)=\sum_{h=1}^{H_{\mathrm{test}}}\log \Pr(Y_h\mid Y_{-h},M_k),
\end{equation}
where $Y_{-h}$ is retained data. Equation~\eqref{eq:lps} favors models that predict omitted genomic blocks, fossils, populations, or laboratories rather than only explaining a selected observation. Holdout units are grouped to respect linkage and kinship; random individual-level splitting can leak closely related information.

For an estimand $\widehat{\psi}$ under assumptions $a$ and $a'$, report a normalized sensitivity score
\begin{equation}
\label{eq:sensitivity}
\Delta(a,a')=\frac{\left|\widehat{\psi}(a)-\widehat{\psi}(a')\right|}{\max\left(\left|\widehat{\psi}(a)\right|,\varepsilon\right)},
\end{equation}
where $a'$ changes a contamination bound, map, mask, source panel, chronology, or demographic prior, and $\varepsilon>0$ prevents a degenerate denominator. Equation~\eqref{eq:sensitivity} identifies which scientific conclusion is stable, not merely which computation is stable. A large value narrows the claim or specifies the next needed sample.

For a binary model prediction $\pi_r$ and held-out observation $y_r$, calibration can be summarized by
\begin{equation}
\label{eq:brier}
\mathrm{BS}=\frac{1}{R}\sum_{r=1}^{R}(y_r-\pi_r)^2.
\end{equation}
Equation~\eqref{eq:brier} is used for model-class predictions or simulation-class recovery, not to reduce complex ancestry to a single binary trait. The report also gives interval coverage, residual patterns by coverage and chronology, and model disagreement.

\begin{table*}[t]
\caption{Outcome-to-interpretation matrix for ARCHAIC.}
\label{tab:outcomes}
\centering
\footnotesize
\begin{tabular}{L{0.27\textwidth} L{0.27\textwidth} L{0.31\textwidth}}
\toprule
Observed outcome & Permitted interpretation & Required next step or restriction \\
\midrule
Data fail authenticity or provenance gate & no lineage inference from that sample & improve authentication, obtain permitted replicate material, or exclude according to preregistered rule \\
D/$f_4$ asymmetry without demographic comparison & descriptive departure from specified tree expectation & compare ILS, ancestral structure, and alternative gene-flow models \\
Gene-flow model predicts held-out data better than alternatives & gene-flow-supported under the declared model set & do not claim an exact fossil donor or date without local/temporal evidence \\
Tract/local genealogy inference is stable across methods and maps & source-qualified/time-supported conditional result & report intervals, source-family uncertainty, and pulse-model dependence \\
Recent chromosome-scale ancestry in authenticated individual & direct-hybrid evidence for that individual & do not generalize to population-wide mating rates without population sampling \\
Introgressed region with selection or function signal & adaptive-retention candidate only & replicate and separate demography, linked selection, and phenotype ascertainment \\
\bottomrule
\end{tabular}
\end{table*}

\subsection{Functional consequences are a separate module}

Introgression can introduce variants that later drift, decline under purifying selection, or increase under selection. The genomic landscape of Neanderthal ancestry is nonuniform, and both selection and demographic parameters affect its distribution~\cite{sankararaman2014,juric2016}. Therefore a functional module starts after, not before, ancestry inference. It prespecifies annotations, population comparisons, gene-set hypotheses, covariates, phenotype ascertainment, and a replication cohort. A molecular annotation is not a clinical, cognitive, or social interpretation.

The EPAS1 analysis shows how a tract can motivate a biological hypothesis when source similarity, population frequency, and phenotype-relevant selection are jointly considered~\cite{huerta2014}. The appropriate conclusion is still conditional: source attribution, selection, regulatory mechanism, and organismal effect each carry different uncertainties. The present design does not measure any of them.

\section{Governance, Reproducibility, and Discussion}

Ancient human remains are not ordinary archival samples. Destructive sampling must be justified by expected information gain, minimal material use, curatorial approval, local legal requirements, and engagement with relevant communities and descendant groups where applicable. Present-day genome panels require consent, controlled access, and careful communication. Ancestry estimates should not be used to attach biological determinism, hierarchy, or social value to present-day communities.

Reproducibility has four layers: physical sample and laboratory record; curated sequence and genotype-likelihood data; locked software, reference, mask, and model specification; and an evidence-labeled report. Independent validation can mean a second library, alternate mapping/reference, a different ancient specimen, another laboratory, a new population, or a held-out genomic unit. It must state what has and has not been replicated. Negative results, failed local-ancestry agreement, and source ambiguity are released in the same ledger as positive results where legal and ethical constraints permit.

\begin{table*}[t]
\caption{Failure-oriented decision schedule for an ARCHAIC program.}
\label{tab:schedule}
\centering
\footnotesize
\begin{tabular}{L{0.18\textwidth} L{0.24\textwidth} L{0.27\textwidth} L{0.21\textwidth}}
\toprule
Module & Decision question & Evidence required for advancement & Action when criterion fails \\
\midrule
Authentication & Is endogenous ancient sequence sufficiently bounded? & age/provenance record, damage/contamination analyses, coverage and callability & retain only a material-quality report; do not fit ancestry model \\
Genome-wide signal & Does a sharing statistic exceed tree expectation robustly? & prespecified quartet, blocks, filters, and error sensitivity & retain descriptive statistic; expand alternative demographic models \\
Demography & Does a gene-flow model predict better than structure/ILS alternatives? & simulation and held-out support under explicit model set & downgrade to unresolved history; collect discriminating samples \\
Local ancestry/timing & Are tracts and dates stable across methods/maps? & agreement, source sensitivity, recombination and chronology intervals & use source-family and broad-time language only \\
Direct hybrid & Does individual data support recent parentage? & chromosome-scale likelihood ratio against matched artifacts & report population-level ancestry only if supported \\
Function & Is retained ancestry linked to selection/function beyond confounding? & separate preregistered selection, annotation, and replication evidence & prohibit adaptive or phenotype claim \\
Independent evaluation & Does another workflow reproduce the label? & new sample, library, method, or analyst reproduces relevant gate & mark inference provisional and locate disagreement \\
\bottomrule
\end{tabular}
\end{table*}

The principal scientific contribution of this framework is a disciplined bridge from a high-level historical statement to its observable implications. A direct hybrid, D statistic, f4-ratio, HMM tract, ARG local tree, tract date, and functional assay are not interchangeable. Each can be informative while leaving another link unresolved. The strength of ancient DNA is that dated genomes can discriminate models that present-day data leave entangled; its limit is that preservation, sampling, and source coverage remain uneven. \archaic{} makes those limits visible rather than hiding them behind a final admixture diagram.

\section{Conclusion}

Ancient humans did interbreed with other hominins. Genomic evidence supports gene flow between Neanderthal-related and modern-human-related populations, Denisovan-related and modern-human-related populations, and direct Neanderthal--Denisovan parentage for Denisova 11. The scientifically responsible next question is not whether to repeat that affirmative answer, but how to identify lineage pair, direction, timing, pulse structure, and retained effects without conflating them. \archaic{} provides a staged route: authenticate material, define references, compare demographic alternatives, reconstruct local ancestry, estimate conditional direction and timing, validate independently, and calibrate the claim. It does not deliver new ancestry results; it defines how future research can do so transparently.

\appendices
\section{Minimum Preregistration and Reporting Record}

Before any primary statistic is computed, the study registers the scientific target: focal recipient and source families, individual/site/time-bin distinction, outgroup, geographic and calendar-time scope, genomic build, callable-region mask, and intended claim label. It records all specimens considered, the reason for inclusion or exclusion, provenance and age metadata, extraction/library identifier, coverage, damage and contamination procedures, sex-chromosome policy, and whether observations are genotype likelihoods or pseudo-haploid calls. This record prevents a sample, label, or mask from being silently changed after a salient result is seen.

For genome-wide inference, the ledger specifies quartets, block size, primary D and $f_4$ calculations, source graph, tree/structure/migration alternatives, simulation design, priors or parameter ranges, and the unit held out for prediction. For local ancestry, it names the reference panels, algorithms, recombination maps, tract filters, and agreement rule. For chronology, it records the generation-time and map assumptions, any ancient-sample temporal anchors, and the condition under which only a broad interval or source-family label is allowed. For a direct-hybrid analysis, it fixes the simulated parentage, ILS, and contamination alternatives before inspecting chromosome-scale results.

\subsection{Deviation and evidence ledger}

A deviation record states its date, trigger, affected \archaic{} stage, changed sample, mask, reference, or model, and whether it was made before or after a relevant outcome was examined. For example, changing the source panel after observing a high local-ancestry score is logged as a reference modification; adding a migration edge after a tree model fails is logged as a model-set expansion. The final report provides original and updated specifications and states whether the evidence label changes under both. This distinguishes practical adaptation to a preservation constraint from outcome-driven reinterpretation.

Each material claim receives linked identifiers for the sample/passport, genomic region or statistic, reference panel, comparison model, uncertainty interval, and code/data release version. A gene-flow-supported row points to the ILS and structure alternatives it survived. A source-qualified/time-supported row points to local-ancestry agreement, recombination sensitivity, and dated samples. A functional row points separately to selection, annotation, and replication evidence. Evidence that cannot be re-linked to these identifiers may remain useful background but cannot promote a historical claim.

\subsection{Source uncertainty and controlled release}

The report separates a sampled reference individual from the unknown historical donor population. When a segment is more similar to a Neanderthal or Denisovan reference than to other references, the preferred language is ``Neanderthal-related'' or ``Denisovan-related'' unless the source alternatives have been shown to be identifiable. A source-qualified statement must list the sampled references, the unobserved-source model, the date and recombination assumptions, and the alternative graphs that produce materially similar predictions. If a ghost-source model improves fit but does not identify a unique branch, the output is a model feature rather than a new fossil population or taxonomic label.

Uncertainty is released as an object, not an afterthought. Each principal statistic is accompanied by callable-site counts, block or bootstrap configuration, genotype-likelihood treatment, contamination interval, and the result of transversion, library, mapping, and reference-panel sensitivity checks. Each local-ancestry result carries a posterior or support score, agreement rate across algorithms, and the rule by which uncertain segments were masked. Each timing result includes the recombination maps, generation-time range, pulse model, and interval. Claims are downgraded automatically when an interval overlaps a meaningful demographic alternative or when results change direction under a prespecified quality filter.

Sequence reads from ancient human remains, location metadata, and present-day genome panels can have distinct access restrictions. The reproducibility package therefore releases code, parameter files, masks, simulation seeds, non-sensitive summary statistics, and a machine-readable deviation ledger whenever permitted, while providing a governed access path for restricted data. The publication explains which claims can be independently checked from open summaries and which require authorized data access. This practice permits scientific scrutiny without treating rare remains, sensitive sites, or living communities as unrestricted resources.

\subsection{Simulation recovery suite}

Before empirical interpretation, the analysis runs a simulation recovery suite matched to the planned chronology, coverage distribution, damage profile, contamination bounds, recombination maps, and reference-panel gaps. The suite includes a strict tree with incomplete lineage sorting, a structured ancestral population, single and multiple admixture pulses, reciprocal gene flow, and an unsampled-source scenario. It also varies post-admixture selection, population size, migration, and low-coverage calling. This design asks a practical question: under the data quality actually available, can a procedure distinguish the historical feature it is intended to estimate from a nearby alternative? A method that recovers a pulse only when its source is perfectly sampled cannot justify a strong exact-source claim for ancient hominins.

The recovery report records power, false-positive rate, interval coverage, source-family confusion, direction error, and timing bias for every preregistered data-quality stratum. For local ancestry, it records segment precision and recall conditional on a simulated truth, the fraction of uncertain segments, and sensitivity to the reference panels and recombination maps. For D and $f_4$ statistics, it records calibration under linkage, block definition, genotype uncertainty, ancestral structure, and contamination. For direct-hybrid analysis, it generates first-generation, more distant admixed, ILS-only, and contaminant-mixed genomes at matched coverage so that the likelihood ratio in \eqref{eq:hybridlr} has a known error profile. Simulation is not used to declare a past event true; it states what patterns a method can and cannot resolve.

The release criterion is tied to the claim ladder. If a tree statistic is calibrated but pulse number is not recoverable, the report may state allele-sharing asymmetry or gene-flow support but not a unique pulse history. If a source family is recovered while exact references are frequently confused, the report uses source-family terminology. If dates lose coverage under plausible recombination maps, the output is a broad time interval or no timing claim. If direct-hybrid simulations cannot separate recent parentage from an artifact at the available coverage, the individual is not labeled a direct hybrid. These rules identify whether the next improvement is a new sample, library, source panel, map, or narrower conclusion.

\balance
\begin{thebibliography}{99}

\bibitem{green2010}
R. E. Green \emph{et al.}, ``A draft sequence of the Neandertal genome,'' \emph{Science}, vol. 328, no. 5979, pp. 710--722, 2010, doi: 10.1126/science.1188021.

\bibitem{reich2010}
D. Reich \emph{et al.}, ``Genetic history of an archaic hominin group from Denisova Cave in Siberia,'' \emph{Nature}, vol. 468, pp. 1053--1060, 2010, doi: 10.1038/nature09710.

\bibitem{slon2018}
V. Slon \emph{et al.}, ``The genome of the offspring of a Neanderthal mother and a Denisovan father,'' \emph{Nature}, vol. 561, pp. 113--116, 2018, doi: 10.1038/s41586-018-0455-x.

\bibitem{fu2015}
Q. Fu \emph{et al.}, ``An early modern human from Romania with a recent Neanderthal ancestor,'' \emph{Nature}, vol. 524, pp. 216--219, 2015, doi: 10.1038/nature14558.

\bibitem{iasi2021}
L. N. Iasi \emph{et al.}, ``Initial Upper Palaeolithic humans in Europe had recent Neanderthal ancestry,'' \emph{Nature}, vol. 592, pp. 253--257, 2021, doi: 10.1038/s41586-021-03335-3.

\bibitem{prufer2014}
K. Pr\"ufer \emph{et al.}, ``The complete genome sequence of a Neanderthal from the Altai Mountains,'' \emph{Nature}, vol. 505, pp. 43--49, 2014, doi: 10.1038/nature12886.

\bibitem{meyer2012}
M. Meyer \emph{et al.}, ``A high-coverage genome sequence from an archaic Denisovan individual,'' \emph{Science}, vol. 338, no. 6104, pp. 222--226, 2012, doi: 10.1126/science.1224344.

\bibitem{vernot2014}
B. Vernot and J. M. Akey, ``Resurrecting surviving Neandertal lineages from modern human genomes,'' \emph{Science}, vol. 343, no. 6174, pp. 1017--1021, 2014, doi: 10.1126/science.1245938.

\bibitem{sankararaman2014}
S. Sankararaman, S. Mallick, M. Dannemann, K. Pr\"ufer, J. Kelso, S. P\"a\"abo, N. Patterson, and D. Reich, ``The genomic landscape of Neanderthal ancestry in present-day humans,'' \emph{Nature}, vol. 507, pp. 354--357, 2014, doi: 10.1038/nature12961.

\bibitem{juric2016}
I. Juri\'c, S. Aeschbacher, and G. Coop, ``The strength of selection against Neanderthal introgression,'' \emph{PLoS Genetics}, vol. 12, no. 11, e1006340, 2016, doi: 10.1371/journal.pgen.1006340.

\bibitem{bergstrom2020}
A. Bergstr\"om \emph{et al.}, ``Insights into human genetic variation and population history from 929 diverse genomes,'' \emph{Science}, vol. 367, no. 6484, eaay5012, 2020, doi: 10.1126/science.aay5012.

\bibitem{huerta2014}
E. Huerta-S\'anchez \emph{et al.}, ``Altitude adaptation in Tibetans caused by introgression of Denisovan-like DNA,'' \emph{Nature}, vol. 512, pp. 194--197, 2014, doi: 10.1038/nature13408.

\bibitem{kuhlwilm2016}
M. Kuhlwilm \emph{et al.}, ``Ancient gene flow from early modern humans into Eastern Neanderthals,'' \emph{Nature}, vol. 530, pp. 429--433, 2016, doi: 10.1038/nature16544.

\bibitem{hubisz2020}
M. J. Hubisz, A. L. Williams, and A. Siepel, ``Mapping gene flow between ancient hominins through demography-aware inference of the ancestral recombination graph,'' \emph{PLoS Genetics}, vol. 16, no. 8, e1008895, 2020, doi: 10.1371/journal.pgen.1008895.

\bibitem{schaefer2021}
N. K. Schaefer, B. Shapiro, and R. E. Green, ``An ancestral recombination graph of human, Neanderthal, and Denisovan genomes,'' \emph{Science Advances}, vol. 7, no. 29, eabc0776, 2021, doi: 10.1126/sciadv.abc0776.

\bibitem{durand2011}
E. Y. Durand, N. Patterson, D. Reich, and M. Slatkin, ``Testing for ancient admixture between closely related populations,'' \emph{Molecular Biology and Evolution}, vol. 28, no. 8, pp. 2239--2252, 2011, doi: 10.1093/molbev/msr048.

\bibitem{patterson2012}
N. Patterson \emph{et al.}, ``Ancient admixture in human history,'' \emph{Genetics}, vol. 192, no. 3, pp. 1065--1093, 2012, doi: 10.1534/genetics.112.145037.

\bibitem{sumer2024}
A. P. S\"umer \emph{et al.}, ``Earliest modern human genomes constrain timing of Neanderthal admixture,'' \emph{Nature}, vol. 638, pp. 711--717, 2024, doi: 10.1038/s41586-024-08420-x.

\end{thebibliography}

\end{document}
